% Integration report for Dot batches 69--87.  See integration/late_integration.json.
\head{第十八版新增续稿译校附记（Dot 69—87）}
\egasection{本次整合与定点复核范围}
本组把已经存在的十九个独立续译包接入合订稿：第69—72批为第二篇附录17.16，第73—85批为第三篇，第86批为1965年9月29日致迪厄多内的信，第87批为邻接手稿及卷首异扫。作者正文、证明、原稿写作意见、可辨残文及原独立译稿的脚注均保留。独立包末的长篇译校说明移入以下对应小节；这些补证是译校文字，并非格罗滕迪克原稿。本次接入不等于对上述全部数学结论完成新的独立审定。

本次实际查看的底本图像为：英译物理第78、80、81、86、87、90、106—109页；法文物理第1、126、136、140、141、143—145、150、153、161、162、164页。以下只记录由这些页面以及对应译稿逐段读取得到的定点结论；其他脚注中的历史核查、外部文献和补证保留其原有来源身份，不把它们全部改称本轮重新验证的结果。

\begin{enumerate}[label=\arabic*.]
\item \textbf{16.2、16.3及16.9。} 英78、法161的16.2原句没有正秩限定；原独立译稿依15.1的定义补加正秩，并以脚注明示，这一处理予以保留。法162的16.3确写作二次型的二变量记号以及秩减一；正文沿用二次型的单变量记号和射影超曲面维数“秩减二”，原读法及几何理由留在脚注。英80—81、法164把恒秩条件放在16.9的第三项附近，整合稿明确先列前两项的等价，再在恒秩条件下加入后两项；没有把所有局部自由模都无条件换成同一秩的标准型。
\item \textbf{16.8的直和系数。} 法164作直和，英80误作张量积；两页在两奇秩情形都给出系数2。原独立稿改为系数4，并给出两个秩一型的直接检验，本次保留其正确正文及原错脚注。16.6(c)的法稿上标痕迹仍按原独立注的谨慎说法记录；没有把该墨迹强行判定为另一条数学断言。
\item \textbf{第三篇的划销稿层。} 英87在第1节开头注明一段在原稿中已被划去，英90在零次态射段末明确结束该范围。原包把起点说明放在封面范围段，整合时将其移到实际正文起点，并保留终点说明；这样既保留材料，也不混淆其稿层。
\item \textbf{命题4.3。} 法140在总假设中手加“局部”，而(a)另加“有限表示”；法141在(b)明写满射。原包据此区分两层假设，予以保留。法144及英109的无穷小论证有一处将“余切映射满射”漏写为原态射满射；原包已补星号。法145和英109的相应 Hom 式把靶写到射影空间的结构层，整合稿沿用原包按拉回后的定义域改成源空间结构层的订正。原作者关于未来无穷小理论及避免循环依赖的写作说明完整保留；本轮没有把这段略证重新宣称为无条件全证。
\item \textbf{书末材料与同文去重。} 法150的草记和倒写笔记仍单列为来源补遗，未辨定词句及公式保留明确标记。法153与法1的题名、目录和末段为同文异扫，并非新的一章；其完整对照译文移到L18-087。第87批原包中的旧“入册方式”“下一处”不作为本版现行进度，完整原件仍保留在归档中。
\end{enumerate}

另按原页新补译三处先前未录的划销材料：法126页中上部从“设$X$在$S$上平坦且局部有限表示”至“至迟在第五章”的整块旧引入；法136页下部从“命题2(b)向我们表明”至“称为与除子线性系相关联的可逆模”的旧论证；法140页中部带原编号5.3的旧推论及其解释。本版在来源补遗之后另设“第三篇删划草稿补译（非定稿）”一章，完整译出这三块可辨主体。三处原独立稿中的未收入说明已改为对应来源引用。法136上半的密集覆打、两个关于$D$的缺损条件、法140斜边注及法126的小笔迹仍保留未辨标记；这个补译范围不等于对法文171页作了全面遗漏审计，也不以英译编辑稿到达末页宣布全部手稿完成。新章中另外标明了参数模的对偶、恢复可逆模前的负扭转以及旧推论5.3弱假设不足三项处理，具体图像位置和依据见工程的划销块记录。

第三篇保留原印的“命题3.1”和另一个“引理3.1”；电子锚点对后者使用单独范围，以避免超链接混淆，不改动纸面编号。各批的特殊数学宏只在本批局部组内定义，不改变合订稿已有宏的意义。独立包封面、重复的“续”标题及人工分页命令不移入正文；这种版面整理不删数学段落。少量“下方补证”“上一批”改为本报告的确切位置；显示公式之后孤立的脚注引标移到公式前的引导句，脚注内容及公式保持不变。完整行号映射、散列、锚点及脚注编号见工程中的整合记录。

\begingroup
\def\OO{\mathcal O}
\def\Sym{\operatorname{Sym}}
\def\rk{\operatorname{rk}}
\egasection{L18-069　极化与判别式：16.1—16.3}
\hypertarget{audit:L18-069}{}
\source{原独立续译第69批；英稿物理页 78,79；法稿物理页 161,162}

\paragraph{极化与判别式。}
此处的约定是
\[
B(v,w)=Q(v+w)-Q(v)-Q(w),
\]
没有因子$1/2$；从而在任意剩余特征中均有
$dQ_x(w)=B(x,w)$。矩阵为该配对在一组坐标中的系数矩阵；如果写成$Q=\sum_i a_i X_i^2+\sum_{i<j} b_{ij}X_iX_j$，则其对角项为$2a_i$，非对角项为$b_{ij}$（按对称方式排列）。这里判别式取极化矩阵的行列式。原稿的保留语“如果没有弄错的话”附于名称；英译将它印成脚注72，英译脚注73则来自“立即”一语，二者不是另立的数学命题。

\paragraph{射影切空间如何检验。}
记$W=E^\vee$。在$x'$选取仿射坐标片，用线性函数$\ell$令$\ell(x)=1$（不改变所代表的射影点）。片上$\ell=1$给出线性补空间$\ker\ell$作为切向量空间。因为
\[
B(x,x)=2Q(x)=0,
\]
$B(x,-)$在$kx$上为零，故限制到$\ker\ell$与取商$W/kx$得到同一非零性判据。非零的$Q$在此给出一个超曲面；它在此点光滑且局部维数为$\rk E-2$，恰当且仅当该微分不为零。若$Q=0$，则零概形是整个射影空间，虽光滑但其维数为$\rk E-1$。特征$2$、偶秩且$N\ne0$时，在代数闭域上$Q|_N=\sum a_i X_i^2=(\sum\sqrt{a_i}X_i)^2$；根基维数至少为$2$，这个线性型必有一个非零零点。不能由此断言原域上一定有这样的非零零点。

\paragraph{相对基底。}
按15.1的用语，正秩时先要求$Q$在每一纤维上非零，再要求射影零概形光滑；秩一非零型的空零概形也光滑。每个纤维中非零方程所定义的射影零概形是有效 Cartier 除子（允许为空），于是相对零概形局部有限表示且平坦，见 \href{https://stacks.math.columbia.edu/tag/062Y}{相对有效 Cartier 除子判据}。在这个前提下，再由\href{https://stacks.math.columbia.edu/tag/01V8}{平坦有限表示态射的逐纤维光滑判据}从几何纤维的论证得出相对结论。非退化时，正秩保证各纤维的$Q$不为零。
\par\endgroup

\begingroup
\def\OO{\mathcal O}
\def\Sym{\operatorname{Sym}}
\egasection{L18-070　校正判别式：16.4—16.8}
\hypertarget{audit:L18-070}{}
\source{原独立续译第70批；英稿物理页 79,80；法稿物理页 162,163,164}

\paragraph{整系数多项式与特征2。}
令$R=\mathbb Z[a_i,b_{ij}]$。极化矩阵$B$的对角线为$2a_i$，对称的非对角元为$b_{ij}$。$R/2R$是整环；其分式域中的$B$是奇阶交替矩阵，秩为偶数，故$\det B=0$，再由单射回到$R/2R$。因此$\det B$的每个整数系数都可除以2。先在$R$中除系数，再代入目标环；这不是在特征2的环中对$d(q)$的值做除法。校正值在特征2仍可能非零。比如
\[
\begin{gathered}
Q=aX^2+bY^2+cZ^2+dXY+eXZ+fYZ\\
\quad\Longrightarrow\quad
\widetilde d(Q)=4abc-af^2-be^2-cd^2+def .
\end{gathered}
\]
模2后，右侧为$Q(f,e,d)$。取$Q=X^2+YZ$，整系数普通、校正判别式分别为$-2,-1$，到特征2后分别为$0,1$。现代术语参见 \href{https://doi.org/10.2140/ant.2018.12.99}{Dolgachev--Duncan，命题2.1--2.2、注2.3--2.4及例2.5}；该文明确引用本附录，称特征2的奇秩不变量为\emph{半判别式}，并记录\emph{半行列式}这一名称。

\paragraph{从普遍矩阵到线丛截面。}
以$G$表示坐标变换，$B\mapsto G^{\mathsf T}BG$，故$d\mapsto(\det G)^2d$。奇秩时，在无整数挠的普遍换基环$R[G_{ij},(\det G)^{-1}]$中消去因子2，之后再特化，由此得到$\widetilde d$的同一换基律。因此局部计算确实粘合为$(\det E)^{\otimes2}$的截面。唯一性也先在无整数挠的普遍系数环中由$2\widetilde d=d$得到，再由基变换推广；在一个固定的含2挠环中，仅有该等式并不决定$\widetilde d$。

\paragraph{直和与系数。}
普通判别式在直和下相乘。偶与偶相加，两侧均为普通判别式的乘积；一奇一偶相加，在普遍系数环中从两侧消去一个2，再特化；奇与奇相加，左侧为普通判别式，因而两个因子2都保留。三种情况均适用于任意基底，不要求2可逆。
\par\endgroup

\begingroup
\def\OO{\mathcal O}
\def\Qsp{\mathcal Q}
\def\Sym{\operatorname{Sym}}
\def\Isom{\operatorname{Isom}}
\def\Hom{\operatorname{Hom}}
\egasection{L18-071　光滑二次型判据：16.9—16.11.2}
\hypertarget{audit:L18-071}{}
\source{原独立续译第71批；英稿物理页 80,81,82,83；法稿物理页 164,165,166,167,168}

\paragraph{同构空间及相对基底。}
局部平凡化后，在矩阵坐标环中使行列式可逆（可添变量$t$并加$t\det(u)=1$），再施加两型相容的有限个系数方程，得到仿射有限表示的$P$。秩不恒定时，(i)及(ii)按各开闭秩层分别检查。正秩时，二次型的光滑条件可在几何纤维上检验（参看本版核查报告 L18-069 的相对 Cartier 平坦性说明）；线丛截面可逆也等价于在每个几何纤维上非零。对(iii)所用的满射不必平坦：在每点取一上方点，截面若原来在纤维上为零，随剩余域扩张仍为零。推论16.10的有限扩域可由非空有限型同构预概形的一个闭点取得；不声称扩域可分。

\paragraph{轨道映射的平坦和满射。}
两端均为向量丛的开子概形，在$S$上光滑。按\href{https://stacks.math.columbia.edu/tag/039C}{有限表示的逐纤维平坦性判据}，只需在几何纤维证明平坦。在代数闭域上，目标是整的仿射空间的非空开子概形，故可用泛平坦性取$U$。等变性使所有$u(U)$之上平坦。$U$含$k$点，传递性使这些开集的并包含所有闭点；非空闭补集必有闭点，故确实覆盖。满射性可在任一点的代数闭剩余域上用16.11.1，获得其上方的一点。

\paragraph{奇秩根基。}
在特征2的代数闭域上，$Q|_N$是一个线性型的平方；若$\dim N\ge2$，必有非零零点。奇秩交替型的根基非零，所以光滑时恰为一维。取补空间后得到非退化的偶秩极化型，与根基正交。秩一部分重缩放即可化成$X^2$；补空间为零时没有配对项。通常的非退化型分类也可直接得到：维数至少为二时，在二平面上限制齐次二次式，可找非零零向量$x$。选$y$使$B(x,y)=1$，再把$y$换成$y-Q(y)x$，两向量即为一个双曲对。取正交补重复此过程；若最终余一维，非零平方系数可开平方而归一。

\paragraph{fppf 不等于光滑。}
特征2、奇秩时，取标准型的非零极化根基向量$z$。在单位元以双数写$g=1+\varepsilon A$、$\varepsilon^2=0$，轨道映射的一阶变化（按预合成约定）为$v\mapsto B_0(v,Av)$，在$z$恒为零。目标是二次型空间的开集；取$\ell(z)=1$，则切方向$\ell^2$不在像中。因此此态射虽忠实平坦，却不一定光滑；秩一情形就是$\mathbb G_m$上的平方映射。这不妨碍$Q_0$作为二次型光滑。
\par\endgroup

\begingroup
\def\OO{\mathcal O}
\def\Qsp{\mathcal Q}
\def\Sym{\operatorname{Sym}}
\def\Isom{\operatorname{Isom}}
\def\Spec{\operatorname{Spec}}
\egasection{L18-072　正交群、挠子与局部等平凡性：16.12—16.14}
\hypertarget{audit:L18-072}{}
\source{原独立续译第72批；英稿物理页 83,84,85；法稿物理页 168,169,170,171}
\noindent\textbf{一、术语与原文的开放问题。}
作者希望用“判别式”指刻画光滑性的校正量，将极化双线性型的行列式另称“行列式”。现代术语并不完全统一：Dolgachev--Duncan 对奇维特征 $2$ 使用 \emph{half-discriminant}，并回顾相关术语和归功。此旁证不代替作者的历史提问，也不把不同归一化下的判别式直接视为同一个量。参见 \href{https://doi.org/10.2140/ant.2018.12.99}{Dolgachev--Duncan，命题2.1--2.2及评注2.3--2.4}。

\medskip
\noindent\textbf{二、16.12 的挠子判据。}
令 $I=\Isom(E_0,E)$，$Y=\Qsp(E)^*$。右作用给出
\[
 I\times_S \mathrm O(n)_S\longrightarrow I\times_Y I,
 \qquad (u,v)\longmapsto(u,uv).
\]
对任一测试概形，逆映射是 $(u,u')\mapsto(u,u^{-1}u')$：两者输运的标准型相同，恰意味着 $u^{-1}u'$ 保持标准型；等式不只在域值点成立。逆映射自然地与拉回相容，故给出预概形同构。再由16.11，$I\to Y$ 是 fppf。沿 $Q$ 的截面拉回即得 $P$。反过来，沿右挠子下降标准二次空间；用标准左表示形成伴随空间时，等价关系写作 $(pv,x)\sim(p,vx)$，因此等距同构挠子的评价映射 $[p,x]\mapsto p(x)$ 定义良好。fppf 下降给出互逆构造及其基变换相容性。

\medskip

\noindent\textbf{三、平展与有限平展不是同一个断言。}
光滑满射可在平展局部取截面；有限平展平凡化另用 \href{https://grothendiecksga.com/read/sga3/fr/XXIV-4.html}{SGA 3，第二十四讲，4.2.2--4.2.4}。此处标准 $\mathrm{SO}(n)$ 在 $n\ge3$ 时为定型半单群，$n=2$ 时为分裂环面，$n=1$ 时平凡；后两种同样可用分裂约化群情形（4.2.3）。在16.13的假设下，商是有限平展群，先拉回商挠子，再应用该定理；局部化后所得有限平展覆盖仍半局部。对 Zariski 局部陈述，或按4.1.2等价判据，或利用有限表示将局部环上的数据扩到某个开邻域。第16.14项在奇秩特征 $2$ 时第一步并不平展；但有限局部自由覆盖后可再作有限平展平凡化，复合仍有限、忠实平坦且有限表示。

\medskip
\noindent\textbf{四、秩一反例与结构群之商。}
秩一的射影空间是 $\mathbb P^0_k$，$at^2$ 的射影零概形为空，故该型按定义光滑；要化为 $t^2$，等价于让 $a$ 成为平方。扩张 $k(\sqrt a)/k$ 是二次纯不可分扩张。它提供 fppf 覆盖，却不能用可分扩张代替。
\par
奇秩时，行列式把正交群映到 $\mu_2$；标量矩阵 $\epsilon I$ 的行列式等于 $\epsilon^n=\epsilon$，故给出截面。局部取 $\det E$ 的平凡化后，若 $u$ 输运标准型为 $Q$，则 $(\det u)^2=\widetilde d(Q)/\widetilde d(Q_0)$。因此局部商挠子正是取该单位的平方根。偶秩时，商由 Dickson 不变量给出，并不在特征 $2$ 处退化为 $\mu_2$。例如二维双曲型 $xy$ 的变量交换在特征 $2$ 处行列式为 $1$，其 Dickson 不变量却非零。参见 \href{https://math.stanford.edu/~conrad/papers/luminysga3.pdf}{Conrad，附录 C.2.8、C.2.11--C.2.12}。这也解释16.13的例外为何恰为奇秩且剩余特征为 $2$。
\par\endgroup

\begingroup
\def\OO{\mathcal O}
\def\LL{\mathcal L}
\def\MM{\mathcal M}
\def\EE{\mathcal E}
\def\Pic{\operatorname{Pic}}
\def\Sym{\operatorname{Sym}}
\def\Proj{\operatorname{Proj}}
\def\Spec{\operatorname{Spec}}
\def\Hom{\operatorname{Hom}}
\egasection{L18-073　第三篇卷首与射影丛的 Picard 群}
\hypertarget{audit:L18-073}{}
\source{原独立续译第73批；英稿物理页 86,87,88,89,90；法稿物理页 151,152,157,154,155,156,116,117,118}
\paragraph{分层与相容同构。}
由1.5知，在秩至少为2的部分，纤维次数局部常值。对每个 \(S_n\) 令 \(P_n=f^{-1}(S_n)\)、\(f_n=f|_{P_n}\)，取
\[
\MM_n=(f_n)_*\bigl(\LL|_{P_n}(-n)\bigr).
\]
由1.5在 \(S_n\) 上的局部平凡性，\(\MM_n\) 可逆；伴随评价映射在这些开集上是同构，所以全局给出 \(f_n^*\MM_n(n)\simeq\LL|_{P_n}\)。这些 \(\MM_n\) 在不交开分割上拼成 \(\MM\)。若只给一个整数 \(n\)，在任一非空、处处正秩的底上，使用同一个 \(f_*(\LL(-n))\) 即可；不能把这个说法解释为可逆模没有标量自同构。底为两部分不交并、两部分上的次数分别为0和1，就是必须使用局部常值次数的例子。

\paragraph{秩为零或一。}
在商约定下，秩一时典范满射 \(f^*\EE\to\OO_P(1)\) 是秩一模之间的同构。识别 \(P=S\) 后，\((*)\) 就是
\(([\MM],n)\mapsto[\MM\otimes\EE^{\otimes n}]\)，核为
\(\{([\EE^{\otimes(-n)}],n):n\in\mathbb Z\}\)。
秩零时 \(P\) 为空，故核是整个源。处处至少一秩时，\(f_*\OO_P\simeq\OO_S\)； 秩为零时不满足此式（除非底也为空）。

\paragraph{域上与局部提升的两个细节。}
当 \(r\ge1\)，删去原点没有删去任何余维一点；在正规仿射空间上，余维一的除子数据决定秩一反身层。由 \(A\) 唯一分解，\(V\) 的任何 Cartier 除子都主，故 \(q^*\LL\) 无分次地平凡；又 \(\Gamma(V,\OO_V)=A\)（在分式域内取余维一局部化之交），所以 \(i_*q^*\LL\) 平凡。若分次 \(A\)-模 \(N\) 无分次地自由秩一，\(N/(t_0,\ldots,t_r)N\) 是一维分次 \(k\)-向量空间。提升它的齐次基元；由于 \(N\) 有限生成且分次下有界，分次 Nakayama 引理说明所提升的元生成 \(N\)，于是它是齐次基。这样的移位恰给 \(\OO(n)\)。

对一般的底，在仿射邻域上可把一个可逆模及其乘法逆的给定同构从 \(\mathbf P_A^r\) 沿有限标准仿射覆盖下降到某个有限生成 \(\mathbb Z\)-子代数：有限表示、有限重叠上的过渡映射及有限等式都只需有限多个系数。诺特情形下，对 \(s\) 的纤维所给的平凡化，在扭掉 \(n\) 后由上同调与基变换及 \(H^1=0\) 可提升为邻域上的截面。它的非生成点构成闭集；因 \(f\) 固有，其像闭且不含 \(s\)。移去该像，截面在整个限制空间上给出平凡化。这说明1.5中“纤维平凡”何以能加强成“在底的邻域上的整个限制空间平凡”。

\paragraph{底本与引用。}
法稿据随工程所附按原文顺序摘出的扫描；英稿见
\href{https://webusers.imj-prg.fr/~leila.schneps/grothendieckcircle/EGA/EGA5trans.pdf#page=86}{第五卷英译编辑稿物理86页起}。
射影丛的商约定和 \(\OO(1)\) 见
\href{https://stacks.math.columbia.edu/tag/01OA}{Stacks Project, Tag 01OA}。
手稿中的 III.4.6.5 与第五章1.1均保留原写法，不把现行编号移入原引文。
\par\endgroup

\begingroup
\def\OO{\mathcal O}
\def\LL{\mathcal L}
\def\MM{\mathcal M}
\def\EE{\mathcal E}
\def\Pic{\operatorname{Pic}}
\def\Sym{\operatorname{Sym}}
\def\Spec{\operatorname{Spec}}
\def\Aut{\operatorname{Aut}}
\def\End{\operatorname{End}}
\def\Div{\operatorname{Div}}
\egasection{L18-074　射影丛同构与相对有效除子：1.6—2.4}
\hypertarget{audit:L18-074}{}
\source{原独立续译第74批；英稿物理页 90,91,92,93,94；法稿物理页 118,119,120,121}
\paragraph{一次与满射。}
对剩余域 \(k\)，若 \(\EE'_k\to\EE_k\otimes\MM_k\) 不满，则其像是一个真线性子空间。有一条非零线性商将这个像杀掉，在该射影点处关联的线性式便不能生成扭转线丛的纤维，矛盾。反过来，满射与典范商的复合仍满；原同态的满射性由所有茎的剩余域和 Nakayama 引理检验。若 \(g\) 已是同构，纤维维数相等，故 \(\EE\) 与 \(\EE'\) 秩在相应层相等；该满射于是为同构。

\paragraph{固定扭转及秩零。}
在有非空纤维的各层可典范地取
\(\MM=f_*(g^*\OO_{P'}(1)\otimes\OO_P(-1))\)。
伴随评价给出所需识别，再推前由 \(g\) 拉回的典范商得到 \(u\)。在秩一层，所有整数都是纤维次数，所以这里必须固定扭转为一次；实际构造在秩一同样给相容的数据。秩零层上的两射影丛均为空，\(g\) 看不见 \(\MM\) 在该层的选择，可取 \(\OO_S\) 在该层的限制。局部平凡化 \(\MM\) 的不同选取只差标量。正秩层上，射影化的自同态空间中由行列式非零定义的标准开集相对仿射，局部就是通常的 \(\mathrm{PGL}\)；局部代表之间相差标量，故与1.8的群函子相合。

\paragraph{纤维正则及参数空间。}
纤维是整射影空间，非零全局截面在其每个仿射平凡化区内给出非零函数，因而不是零因子，并不要求它在纤维上处处不消失。要求的是每个 \(s\) 上的整个多项式 \(\psi(s)\) 不为零。射影丛在底上平坦且局部有限表示；因此，这个纤维条件等价于截面的零位定义一个相对有效 Cartier 除子（参见 \href{https://stacks.math.columbia.edu/tag/062Y}{Stacks Project，062Y}）。有限秩自由坐标中，\(\psi(s)\ne0\) 意味着至少一个系数在 \(s\) 的剩余域非零，因而该系数在邻域上可逆；由此得到局部直和因子和转置满射。次数为零时参数射影丛为 \(S\)，对应空的有效相对除子。对于任意 \(T\to S\)，次数是 \(T\) 上的局部常值函数；所以可表对象须取各次数射影丛的不交和，而不是固定一个全局次数。

\paragraph{引用。}
\href{https://webusers.imj-prg.fr/~leila.schneps/grothendieckcircle/EGA/EGA5trans.pdf#page=90}{英译编辑稿物理90页起}；法文摘页随工程附送。与今用术语、商约定及纤维条件对照，可见
\href{https://stacks.math.columbia.edu/tag/01OA}{Stacks Project：射影丛（01OA）}及
\href{https://stacks.math.columbia.edu/tag/056P}{Stacks Project：相对有效 Cartier 除子（056P）}。
\par\endgroup

\begingroup
\def\OO{\mathcal O}
\def\LL{\mathcal L}
\def\MM{\mathcal M}
\def\EE{\mathcal E}
\def\nn{\boldsymbol n}
\def\Pic{\operatorname{Pic}}
\def\Sym{\operatorname{Sym}}
\def\Div{\operatorname{Div}}
\egasection{L18-075　多重射影丛：2.5—2.10}
\hypertarget{audit:L18-075}{}
\source{原独立续译第75批；英稿物理页 94,95；法稿物理页 121,122,123}
\paragraph{为什么不能丢掉基底因子。}
取 \(S=\mathbf P_k^1\)，取两个秩二平凡因子，令 \(\LL=f^*\OO_S(1)\)。在每个纤维上，\(\LL\) 的多次数都是 \((0,0)\)。但 \(f\) 有截面，沿该截面拉回 \(\LL\) 得到非平凡的 \(\OO_S(1)\)，故 \(\LL\) 不可能等于 \(\OO_P(0,0)\) 的同构类。

\paragraph{如何迭代定理1.1。}
对最后一个因子，先把 \(P\) 看成其余因子乘积 \(B\) 上的射影丛。最后一个次数是 \(B\) 上的局部常值整数函数。由于 \(B\to S\) 开且满射，各纤维非空且连通，这个函数在同一纤维上取常值，并下降成 \(S\) 上的局部常值函数。对其余因子重复。在各纤维上按因子限制即可确定各部分次数，因此次数分割唯一。在每个次数层上消去扭转，再用 \(f_*\OO_P=\OO_S\)，可取
\(\MM=f_*(\LL\otimes\OO_P(-\nn))\)（式中按该层限制 \(f\) 和 \(\LL\)）；伴随评价给出所需的相容同构。

\paragraph{参数空间与非负次数。}
在任一纤维上，若某分量 \(n_i<0\)，对若干因子逐层推前表明全局截面为零，不能定义有效相对除子。对所有 \(n_i\ge0\)，用射影丛的正次推前公式及投影公式逐因子计算，得 \((**)\)，并与基变换相容。各几何纤维都是整的；多齐次截面在每个纤维不恒为零，就在各纤维上正则。射影丛乘积的结构态射平坦、局部有限表示，因此可应用纤维判据。截面等价于从上述对称幂张量积的对偶到 \(\MM\) 的满射；于是得到商约定下的参数空间。零多次数对应空除子。

\paragraph{引用。}
法扫描摘页附在工程内。英译编辑稿见
\href{https://webusers.imj-prg.fr/~leila.schneps/grothendieckcircle/EGA/EGA5trans.pdf#page=94}{物理94页起}。
射影丛商约定和相对有效 Cartier 除子的纤维判据分别见
\href{https://stacks.math.columbia.edu/tag/01OA}{Stacks Project，01OA}、
\href{https://stacks.math.columbia.edu/tag/062Y}{Stacks Project，062Y}。
\par\endgroup

\begingroup
\def\Div{\operatorname{Div}}
\def\Supp{\operatorname{Supp}}
\def\Ass{\operatorname{Ass}}
\def\Spec{\operatorname{Spec}}
\def\PP{\mathbf P}
\def\OO{\mathcal O}
\egasection{L18-078　两个投影与命题3.1}
\hypertarget{audit:L18-078}{}
\source{原独立续译第78批；英稿物理页 98,99；法稿物理页 129,130,131}
按商约定，\(\PP(F)\) 参数化 \(F\) 的可逆商。在 \(P\times_S P^\vee\) 上隐去各个拉回符号，
\(P^\vee=\PP(F)\) 的典范商取对偶，得到线子模 \(\OO_{P^\vee}(-1)\subset F^\vee\)。
把它与 \(P=\PP(F^\vee)\) 的典范商复合：
\[
 \OO_{P^\vee}(-1)\longrightarrow F^\vee
          \longrightarrow\OO_P(1).
\]
这是 \(\OO_P(1)\otimes\OO_{P^\vee}(1)\) 的一个截面，其零概形就是 \(H\)。
固定 \(p\in P\)，该截面在参数射影纤维上是非零线性式；纤维维数为零时其零集为空。
所以 \(H\) 相对于 \(P\) 是相对有效 Cartier 除子。
把它沿 \(f:X\to P\) 基变换，得到的除子相对于 \(X\) 仍是相对有效 Cartier 除子，且没有固定点。

另一方向，固定一个几何参数 \(q\in P^\vee(\bar k)\)。
此时 \(H\) 在 \(P_{\bar k}\) 上的限制是 \(q\) 所表示的超平面。
在 \(X_{\bar k}\) 上，拉回截面的局部方程为非零因子，
当且仅当它不在任一关联点上消失。这又等价于没有关联点的闭包整个映入该超平面。
这一非包含性在代数闭包之后的扩域下保持：关联点的闭包是几何整的，扩域不产生新的关联分支；一个分支上线性式的消失条件是线性方程组，扩域不改变其核是否为零。
由于 \(X\times_S P^\vee\to P^\vee\) 平坦且局部有限表示，
\href{https://stacks.math.columbia.edu/tag/056P}{相对有效 Cartier 除子的纤维判据}
把这个逐纤维条件等价地变为命题中的相对除子条件。
反之，满足此条件的 \(f\) 的拉回给出原侧的除子族；它的转置正是 \(f\) 所分类的超平面族，
所以构造彼此相逆。这也说明若仅检验极小关联点，可能漏掉使方程成为零因子的嵌入关联点。
射影丛的商约定可参看 \href{https://stacks.math.columbia.edu/tag/01OA}{射影丛的定义和典范商}。
\par\endgroup

\begingroup
\def\depth{\operatorname{depth}}
\def\Spec{\operatorname{Spec}}
\def\OO{\mathcal O}
\def\FF{\mathcal F}
\def\JJ{\mathcal J}
\egasection{L18-079　命题3.2与深度引理}
\hypertarget{audit:L18-079}{}
\source{原独立续译第79批；英稿物理页 99,100；法稿物理页 131,132}
原稿的局部化论证不应把任意闭集 \(T\) 不加论证地当成有限表示闭子概形。
下面只下降有限个环元素及有限表示的模，不下降整个 \(T\)。

在 \(x\in T\) 的仿射邻域令 \(I\) 定义 \(T\)，\(s\) 是 \(x\) 的像。纤维的局部环
\(\OO_{X_s,x}\) 诺特。在包含 \(I\) 的该局部环素点处，若纤维模非零，其深度都有假设中的下界；
因此由素避及正则元降深度，能在 \(I\) 的纤维像中找到长度一（或二）的纤维模正则序列。
若 \(\FF_x=0\)，有限生成性使它在某邻域已为零，可直接处理。
若 \(\FF_x\ne0\)，Nakayama 引理保证该点纤维模非零。
清去在该点可逆的分母，可把序列取为 \(I\) 中元素 \(h_1\)（或 \(h_1,h_2\)）的像。

对这些有限数据作诺特逼近：有限表示代数、有限表示模、底上的平坦性及这几个元素
在足够后的仿射阶段下降，见 \href{https://stacks.math.columbia.edu/tag/00R3}{有限表示模的平坦性下降}。
从较早阶段相应纤维的局部环到现有纤维局部环是忠实平坦同态，故正则性在该点下降。
在诺特阶段，\href{https://stacks.math.columbia.edu/tag/00ME}{纤维注入判据}
把逐次乘 \(h_j\) 的纤维注入提升为局部注入，并使逐次商模在底上平坦。
因诺特阶段核有限生成且商模的平坦处为开，缩小到该点的邻域后可同时实现这些断言；
再拉回，商模的平坦性保证短正合列仍正合。
所以在当前底上，局部也可找到这样的模正则序列。

长度一时，\(h_1\) 在模上单射，且主开集 \(D(h_1)\) 在 \(X\setminus T\) 内，
从而限制映射是单射。
长度二时，模 \(M\) 在 \(D(h_1)\cup D(h_2)\) 上的截面由 \(M\) 唯一延拓：
若两个分式 \(a/h_1^m,\ b/h_2^n\) 在交集相等，清去额外分母后可令 \(h_2^n a=h_1^m b\)。
由于 \(h_2\) 在 \(M/h_1^mM\) 上仍为非零因子，\(a=h_1^m c\)；
再由 \(h_1\) 的正则性得 \(b=h_2^n c\)。
因此先把原截面在这两个主开集上的限制延拓，再用 \(h_1\) 在模上单射，
可知该延拓在整个 \(X\setminus T\) 与原截面相同（均在所取的仿射邻域内）。
这就证明了所需的层同构，不需要把 \(X\setminus T\) 假定为拟紧。
\par\endgroup

\begingroup
\def\depth{\operatorname{depth}}
\def\Spec{\operatorname{Spec}}
\def\Ass{\operatorname{Ass}}
\def\OO{\mathcal O}
\def\FF{\mathcal F}
\def\JJ{\mathcal J}
\egasection{L18-080　相对伪态射与可逆层延拓：3.3—3.4}
\hypertarget{audit:L18-080}{}
\source{原独立续译第80批；英稿物理页 100,101,102；法稿物理页 133,134,135}
\paragraph{可逆层、截面和相对除子。}
深度引理给出 \(i_*\OO_U=\OO_X\)，对有限型局部自由层同理；对两层之间的同态应用伴随即得全忠实。任一点的开邻域都与 \(U\) 相交，因为每个纤维的关联点都在 \(U\)；局部常值的秩因此由限制决定。若 \(i_*L'\) 局部自由，则 \(i^*i_*L'\simeq L'\)，故得到 b)。
在 c) 中，若 \(L'\) 可逆延拓为 \(L\)，定义 \(D'\) 的正则截面通过 \(L=i_*L'\) 唯一延拓。乘以该截面给出的 \(L^{-1}\to\OO_X\) 为单射：在 \(U\) 上是单射，推前仍左正合，再用两个层的深度延拓同构。
其像是可逆理想；而且由推前原短正合列的左正合性，它是 \(\OO_X\to i_*\OO_{D'}\) 的核，故定义的正是概形论闭包。若 \(D'\) 是相对有效除子，延拓截面在纤维的所有关联点处是非零因子，故在整个纤维上是非零因子。由平坦且局部有限表示的环境中的纤维判据得到相对有效性；逆向由相对有效除子的基变换稳定性。参见 \href{https://stacks.math.columbia.edu/tag/062Y}{Stacks，062Y}。

\paragraph{不经一般推前基变换的证明。}
给定 \(i_*L_U\) 可逆，令 \(L=i_*L_U\)。
将 \(L\) 拉回到 \(X\times_S Q\)，再张量参数侧 \(\OO_Q(1)\) 的拉回，便直接得到 \(N'\) 的可逆延拓。
纤维深度下界沿到 \(Q\) 后的纤维域扩张保持，故前段构造唯一相对有效除子 \(D\)。其参数侧度数是一：在每个 \(x\in U\) 如此，在全体 \(X\) 各非空正秩层上由度数局部常值及 \(U\) 的稠密性得到（秩一时沿用命题3.1的度一型约定）。固定点不在 \(U\)，故都满足指定深度条件。
反过来，若有 \(D\)，其 \(N=\OO(D)\) 具有参数侧的度一型。在 \(Q\to S\) 的局部平凡区间选取一个截面，使沿 \(X\) 拉回的 \(N\) 再张量参数侧 \(\OO_Q(-1)\) 的相应因子，便在各区间上得到 \(L_U\) 的可逆延拓。深度引理给出这些局部延拓均等于相应的 \(i_*L_U\)，因此 \(i_*L_U\) 可逆。这个论证不要求 \(i\) 拟紧。
\par\endgroup

\begingroup
\def\depth{\operatorname{depth}}
\def\Spec{\operatorname{Spec}}
\def\Ass{\operatorname{Ass}}
\def\OO{\mathcal O}
\def\FF{\mathcal F}
\def\JJ{\mathcal J}
\egasection{L18-081　局部阶乘情形与线性系分类：3.5—4.1}
\hypertarget{audit:L18-081}{}
\source{原独立续译第81批；英稿物理页 102,103,104,105；法稿物理页 135,136,137,138}
\paragraph{局部阶乘与非退化。}
在局部诺特、局部阶乘情形，正规局部整环保证不可约分支局部不相交，可以逐个整分支处理。
由深度条件，\(U\) 含每个分支的泛点和所有余维一的点：被删处的纤维深度至少为二，而本版核查报告 L18-079 的深度比较给出总空间深度也至少为二。
在一个整分支上，可逆层取非零有理截面而写成 Cartier 除子；在 \(U\) 的余维一点读其 Weil 系数，沿余维一闭包原样延拓。所得和在 \(X\) 局部有限：在任一诺特仿射开内，原有理截面在其与 \(U\) 的交上用有限多个仿射平凡片写出，只有有限多个余维一分支。这些 Weil 数据因局部阶乘而给出 Cartier 除子，其线丛限制回原层。
由推论3.4，所得延拓与 \(i_*L_U\) 同构。非退化条件的反例：在代数闭域上把两份不交的仿射直线分别映成 \(\mathbf P^1\) 的两个坐标点。总像不含在单一超平面中，但其中一个超平面的拉回在一整分支上恒零，故不是该纤维的有效 Cartier 除子。

\paragraph{如何检验几何纤维。}
若在一个代数闭包上所有非零向量给出正则截面，则对该闭包上的每个关联点，那个点所处分支上的“给出恒零截面”的向量空间核为零；它可由该分支的一个仿射邻域检验，取局部环模关联素理想即可。
对任意进一步扩域，所考察的有限维向量空间到该整商的映射保持单射。代数闭域上的整有限型仿射分支几何整，扩域后的关联点由原分支及其嵌入关联点的扩张给出，故正则性仍成立；对一般 \(K\) 可再扩到同时包含 \(K\) 和该代数闭包的域，然后按忠实平坦性下降。
几何整情形的简化也可直接看出：若 \(E(s)\to H^0(X_s,L_s)\) 单射，则所有仿射开上取限制的核之交为零。输入有限维，故从这些开取有限个，仍使核之交为零。把到这些仿射开上截面空间的合映射作扩域，保持单射，足以推出扩域后的全局映射单射。反向可由忠实扩域或直接取 \(K=k(s)\) 得到。这个论证不要求 \(X_s\) 整体拟紧。
\par\endgroup

\begingroup
\def\Spec{\operatorname{Spec}}
\def\Hom{\operatorname{Hom}}
\def\Grass{\operatorname{Grass}}
\def\Div{\operatorname{Div}}
\def\OO{\mathcal O}
\def\FF{\mathcal F}
\def\HH{\mathcal{H}\!om}
\def\PP{\mathbf P}
\egasection{L18-082　线性系与格拉斯曼概形：4.2}
\hypertarget{audit:L18-082}{}
\source{原独立续译第82批；英稿物理页 105,106；法稿物理页 138,139,140}
\paragraph{表出模与基变换。}
在仿射基底上，记 \(S=\Spec A\)。在有限表示、固有、平坦的假设下，可把 \(g\) 和可逆层 \(L\) 下降到一个有限生成的 \(\mathbb Z\)-代数上的模型；取足够大的有限阶段，使模型仍固有、平坦且层仍可逆。可从 EGA III\,7.7.5--7.7.6 的有限复形构造作如下理解：对这个诺特模型，取有限投射模组成的非负复形 \(K^\bullet\)，使任意系数模 \(M\) 的零阶截面由 \(H^0(K^\bullet\otimes M)\) 给出。此处使用 \(L\) 对底平坦。这个复形拉回到 \(A\) 后有同样的描述；可以用有限仿射覆盖的平坦 \v{C}ech 复形理解任意换底及任意系数。在仿射基底上便可取
\[
 Q_A=\operatorname{coker}\bigl((K^1)^\vee\longrightarrow (K^0)^\vee\bigr).
\]
事实上，\(K^0,K^1\) 是有限投射模，故 \(\Hom_A(Q_A,M)\) 是
\(\Hom_A((K^0)^\vee,M)\to\Hom_A((K^1)^\vee,M)\) 的核，亦即 \(H^0(K^\bullet\otimes_A M)\)。余核与拉回相容，此构造在任意换底后仍表出相应的系数函子。各仿射开上的表出模因泛性质而唯一同构，满足粘合条件。此处的基变换陈述针对 \(Q\) 和它表出的函子；它并不声称 \(\HH(Q,\OO_S)\) 的形成与任意基变换交换。

\paragraph{从参数丛同构恢复商。}
不同商 \(Q\twoheadrightarrow V\) 给出的射影参数丛图式如果在 \(\PP_S(Q)\) 上同构，从 \(\PP_S(Q)\) 拉回的 \(\OO(1)\) 也随之同构。将该线丛推前到 \(S\)，按射影丛公式恢复 \(V\)，同时将通用商恢复为原来的 \(Q\twoheadrightarrow V\)。这在秩一时同样成立。于是参数丛图式的同构类恰是商模在保持商映射的同构下的等价类，而不是仅按射影丛抽象同构的等价类。

\paragraph{格拉斯曼概形的局部构造。}
在任一仿射开上，取有限表示 \(\OO_S^m\to\OO_S^n\to Q\to0\)。标准秩 \(r>0\) 格拉斯曼概形的通用商，只有在 \(\OO_S^m\) 的像在该商中为零时才经 \(Q\) 分解。这个条件是在通用局部自由商中让有限多个截面同时为零，是闭条件。所得概形由其泛性质粘合，得到 \(\Grass_r(Q)\)。其中 \(r\) 可随底的连通部分局部常值，因而各秩取不交并。任意基变换保持这套陈述；参见 \href{https://stacks.math.columbia.edu/tag/089R}{Stacks，089R} 的标准格拉斯曼概形构造。
\par\endgroup

\begingroup
\def\Spec{\operatorname{Spec}}
\def\Aut{\operatorname{Aut}}
\def\Div{\operatorname{Div}}
\def\OO{\mathcal O}
\def\PP{\mathbf P}
\egasection{L18-083　命题4.3的域底情形与射影标架}
\hypertarget{audit:L18-083}{}
\source{原独立续译第83批；英稿物理页 106,107,108；法稿物理页 140,141,142,143}
\paragraph{限制与底上的单位。}
\(T\) 是代数闭域上的约化整有限型概形，因此几何整。对任一几何参数，原族的截面在 \(T\) 的关联泛点按原纤维正则性不消失；具体地，在含泛点的仿射邻域 \(\Spec A\) 中，关联素理想 \(\mathfrak p\) 给出嵌入 \(A/\mathfrak p\hookrightarrow A\)。扩域后嵌入仍成立，且 \((A/\mathfrak p)\otimes_k K\) 为整环，故扩域后的该泛点仍为关联点。在 \(T\) 的几何纤维上，线丛的非零截面正则。于是该族在 \(T\) 上仍有相对有效除子；限制参数方向的 \(\OO(1)\) 不变。比较原族中相等的除子，可通过方程限制到 \(T\) 后继续比较。
对 \(k\) 上有限表示的 \(T\)，取有限仿射覆盖和其交的有限仿射覆盖；计算全局函数所用的等化子与任何平坦系数扩张交换。对 \(S'\) 的仿射开 \(\Spec R\)，\(R\) 是平坦 \(k\)-模，故 \(H^0(T_R,\OO)=R\)，单位也随之对应。

\paragraph{从相同除子回到参数。}
在 \(T\) 上，\(k\)-线性评估 \(E\to H^0(T,L|_T)\) 单射，因为非零向量的截面正则。把截面限制到 \(T\) 的若干仿射开。由于 \(E\) 有限维，可选有限多个开，使各限制映射的核之交仍为零。对任一 \(S'\) 的仿射开，将合映射与其坐标环作 \(k\)-张量，仍为单射；它通过 \(u\) 在换底后对参数的评估。若 \(u(\psi)=a\,u(\phi)\) 且 \(a\) 是底上的单位，则 \(\psi=a\phi\)，包括 \(S'\) 非约化的情形。相同 Cartier 除子的局部单位比值唯一，因而粘成整体单位。对非代数闭 \(k\) 用忠实平坦的扩域检测两个参数态射是否相等。

\paragraph{标架在任意换底上检测自同构。}
把域上的标架变到坐标点及点 \([1:\cdots:1]\)，再把这些点拉回为 \(S'\) 上的截面。
一个固定所有标架截面的射影自同构，局部可用可逆矩阵表示（用 \(\OO(1)\) 的拉回与底线丛的局部平凡化）。固定各坐标截面迫使矩阵的全部非对角元为零。再固定最后一个截面，所有对角元相等，故该自同构为恒等。这里相等是环上的等式，不是只查看每个点的剩余域；例如在 \(k[\epsilon]/(\epsilon^2)\) 上，矩阵 \(I+\epsilon E_{12}\) 约化为 \(I\)，对应的射影自同构却不恒等。命题4.3的任意底情形接在正文的域底证明之后；相应补证见本版核查报告 L18-084。
\par\endgroup

\begingroup
\def\Spec{\operatorname{Spec}}
\def\Aut{\operatorname{Aut}}
\def\Div{\operatorname{Div}}
\def\OO{\mathcal O}
\def\PP{\mathbf P}
\def\m{\mathfrak m}
\def\Hom{\operatorname{Hom}}
\egasection{L18-084　命题4.3的一般底证明与无穷小差}
\hypertarget{audit:L18-084}{}
\source{原独立续译第84批；英稿物理页 108,109,110；法稿物理页 143,144,145,146}
\paragraph{从所有阶厚化回到诺特底。}
本命题中的 \(F\) 在底上分离。把其对角线沿 \((u,v):S\to F\times_SF\) 拉回，得到由理想层 \(\mathcal I\) 定义的截面等化子。
若所有阶厚化上相等，则 \(\mathcal I_s\subseteq\bigcap_{n\geq0}\m_s^{n+1}\OO_{S,s}=0\)（\href{https://stacks.math.columbia.edu/tag/00IP}{Stacks，00IP}）。这对每个 \(s\) 都成立，故 \(\mathcal I=0\)。

\paragraph{两个提升的差。}
设 \(n\geq1\)。在两截面具有共同像的仿射邻域，两个环映射在模 \(\m^n\) 后相同，差落在 \(V\)。差对乘积满足导子规则，因为两个差的积落在 \(V^2=0\)，而 \(V\) 上的乘法作用经剩余域。它等价于一个从 \(u_0^*\Omega^1_{F_0/k}\) 到 \(V\) 的线性映射。差为零恰表示两提升相等。复合态射时对导子应用链式法则。

对于相对有效 Cartier 除子，局部取两个正则方程 \(b,b'\)。它们在原厚化上的理想相同，故可用一个提升的单位重新归一化，使两方程在模 \(V\) 后相同。两方程之差落在 \(V\otimes_k\OO_{X_0}\)，取其模约化方程 \(b_0\) 的类。这个类定义了从 \(\mathcal I_{D_0}/\mathcal I_{D_0}^2\) 到 \(V\otimes_k\OO_{D_0}\) 的线性映射：换方程只按方程的单位过渡系数改变，换归一化单位则多出一个 \(b_0\) 的倍数。这里用 \(X\) 的平坦性识别增量理想，并用相对有效除子换底后仍是有效 Cartier 除子（\href{https://stacks.math.columbia.edu/tag/056P}{Stacks，056P}）。如果这个类为零，方程相差的增量是约化方程的倍数，可由一个形如 \(1+c\)、\(c\) 落在增量理想中的单位吸收，故两除子相同。局部映射粘合，得到 \(H^0(D_0,\mathcal N_{D_0/X_0})\otimes_kV\) 中的差。

对于态射 \(X\to P\)，可在同像点附近取靶的仿射邻域，比较两环映射；由同样的导子规则，差粘合为 \(\Hom_{\OO_{X_0}}(f_0^*\Omega^1_{P_0/k},\OO_{X_0})\otimes_kV\) 中的元素。此处不要求该全局 Hom 是有限维的；只要 \(V\) 有限维，取一组基便能把相关张量写成有限直和。在对偶数环上得到的是切空间映射；任意此类 \(V\) 的差映射按同样计算为其与 \(V\) 的张量。检查切空间的单射须用纤维函子在对偶数环上的单射，不能改用域值点。b) 的函子态射此处按4.1解释为 \(u\) 的逆转置与基点补开上 \(f\) 的复合。

\paragraph{a) 的任意底：只用像点的纤维。}
在仿射底上，取有限生成的整数子代数作诺特模型。有限表示的 \(X\)、参数丛、两截面、通用除子及两截面所给除子的相等都可在某一有限阶段实现：可先降下闭浸入、线丛及其截面，再用有限表示态射的极限性质使闭子概形的相等成立。由平坦性下降，可使 \(X\) 在模型底上及通用除子在模型参数丛上平坦（\href{https://stacks.math.columbia.edu/tag/01ZL}{Stacks，01ZL}；\href{https://stacks.math.columbia.edu/tag/081C}{Stacks，081C}）。通用除子的线丛方程也可在有限次放大后保持正则：局部乘法映射的余核已对参数底平坦，且源、靶均平坦。于是其核随底扩张而换底；模型上的核有限生成，在极限为零，有限次放大即可消去。故有模型上的相对有效 Cartier 除子，并可用它定义比较函子。

设 \(t\) 是模型底上一个来自实际底点 \(s\) 的像点。模型比较态射在 \(k(t)\) 上的纤维，经扩域到 \(k(s)\) 后，是已完成域底论证的比较态射。由忠实平坦性，它在 \(k(t)\) 上已是单态射：任取同像的两检验截面，扩域后相等，而射影丛的截面相等可沿这种覆盖检验。这不要求模型的所有纤维都满足原稿关于关联循环的几何条件。仅在 \(t\) 处的诺特局部环上作逐阶归纳，只会用这个闭点的纤维；克鲁尔交定理给出两个模型截面在该局部环上相等。拉回至 \(\OO_{S,s}\)，对每个 \(s\) 应用，得所需的相等。

\paragraph{b) 的任意底：可直接检测标架。}
在秩为 \(r>0\) 的部分，令 \(U=X-Z\)。在 \(U\) 的 \((r+1)\) 重纤维积中，像在 \(P\) 里构成射影标架的条件是开条件；局部由各 \(r\) 列行列式为单位给出。这个开子概形在每条几何纤维均非空，由域底标架引理可知。它平坦、局部有限表示、满射于底，所以给出忠实平坦覆盖。在该覆盖上有通用标架截面；若两自同构与 \(f\) 的复合相等，它们在该标架上的作用相同，因此自同构相等。最后沿覆盖检测相等即可。这是 b) 的单独证明，不能直接用于 a) 的参数截面单态射。
\par\endgroup

\begingroup
\def\Div{\operatorname{div}}
\def\Ens{\operatorname{Ens}}
\def\OO{\mathcal O}
\def\KK{\mathcal K}
\def\PP{\mathbf P}
\egasection{L18-085　命题4.4的底除子修订与补证}
\hypertarget{audit:L18-085}{}
\source{原独立续译第85批；英稿物理页 110,111；法稿物理页 146,147}
\noindent\textbf{1. 从除子到截面。}
有效 Cartier 除子 \(B\) 带有可逆层 \(\OO_X(B)\) 及其典范正则截面 \(s_B\)。
反之，可逆层 \(L\) 的正则截面 \(s\) 给出零除子 \(B\)，以及唯一把 \(s_B\)
送到 \(s\) 的同构 \(\OO_X(B)\simeq L\)。局部若 \(B\) 由非零因子 \(b\)
定义，则 \(\OO_X(B)\) 在全分式层中由 \(1/b\) 生成，\(s_B\) 对应函数 \(1\)。
因此亚纯函数 \(f\) 若是全分式层的单位，则 \(f\) 给出 \(\OO_X(B)\)
的全局正则截面，当且仅当 \(B+\Div(f)\) 有效：局部等价于 \(bf\)
是正则函数且为非零因子。\cmt{L18-085-08}{现代记法及正则截面与有效 Cartier 除子的对应见
\href{https://stacks.math.columbia.edu/tag/0C4S}{Stacks Project，0C4S}。
这里“截面正则”指乘法给出单射，不意味着截面处处不消失。}

\medskip\noindent\textbf{2. 证明 a)。}
由4.1，正秩线性系可由 \((L,V)\) 表示，其中 \(V\subseteq H^0(X,L)\)
有限维，且在几何基变换后，\(V\) 的每个非零向量仍给正则截面。参数是 \(V\)
中一维子空间；这是原稿的对偶射影丛记法。若两个除子集合相同，取一个公共除子
\(B\)，选取产生它的截面 \(s\in V\)、\(s'\in V'\)，并用第一步把两个可逆层都
识别为 \(\OO_X(B)\)，使所选截面都对应 \(s_B\)。在同一可逆层上，产生同一
除子的两个正则截面局部相差单位。这些单位唯一，故粘为全局单位，在本题中
属 \(k^\times\)。于是任取 \(0\ne v\in V\)，可从 \(V'\) 选出产生同一除子的
截面，该截面是 \(v\) 的非零标量倍数，故 \(v\in V'\)；交换二者得 \(V=V'\)。
由4.1恢复的参数空间及通用除子相同，故线性系同构。反向由同构对
\(k\)-值点给出双射。若容许空参数，\(k\) 上任意正秩射影空间均有 \(k\)-点，
故两个空集合只对应同一个空参数情形。

\medskip\noindent\textbf{3. 证明 b)。}
给定线性系，选一非零参数截面 \(s_A\)，以其除子为 \(A\)，并把可逆层识别为
\(\OO_X(A)\)。每个非零截面除以 \(s_A\) 得到全分式层的单位；它们与零合成一个
有限维向量子空间 \(E\)，包含 \(1\)。截面的除子正是 \(A+\Div(f)\)，所以
得到右侧的表示。

反过来，由第一步，\(E\) 是 \(H^0(X,\OO_X(A))\) 的有限维子空间，且其中
非零元素给正则截面。若 \(k\) 代数闭，这已是在4.1所需的代数闭包上的检验。
若 \(X\) 几何整，则 \(X_{\overline k}\) 为整概形，扩域后的任何非零向量仍给
非零截面，因而正则。这里用到
\[
 H^0(X,L)\otimes_k\overline k\simeq H^0(X_{\overline k},L_{\overline k}).
\]
有限型的域上预概形拟紧拟分离，取有限仿射覆盖，并用有限仿射开覆盖重叠部分，
可把全局截面写成两条有限乘积间映射的等化子；平坦扩域与这一等化子交换。
于是4.1给出所需的线性系。若无 b) 开头的二选一条件，同一证明仍成立，
但应直接要求几何基变换后所有非零向量给正则截面。

\medskip\noindent\textbf{4. 证明 c)。}
若 \(\Delta=\Delta'\)，取公共除子 \(B=A+\Div(f_0)=A'+\Div(f'_0)\)。
把两套数据都改用 \(B\) 表示，所用函数空间分别为 \(E/f_0\) 与 \(E'/f'_0\)。
它们是在同一可逆层 \(\OO_X(B)\) 中产生同一除子集合的两个截面空间，
故由 a) 的逐截面比较相等。取 \(h=f'_0/f_0\) 即得
\(E'=hE\) 且 \(A'=A-\Div(h)\)。反过来这两式使每个
\(A'+\Div(hf)\) 恰等于 \(A+\Div(f)\)，故集合相同。
若 \(h_1,h_2\) 都给出这两式，\(\Div(h_2/h_1)=0\)，所以
\(h_2/h_1\) 是全局单位，即属 \(k^\times\)；乘任何非零常数也不改变这两式。
\par\endgroup

\begingroup
\def\illis{〔字迹不清〕}
\def\gap{〔扫本缺损〕}
\egasection{L18-087　法文第150页提纲与第153页卷首异扫}
\hypertarget{audit:L18-087}{}
\source{原独立续译第87批；英稿物理页 2,3(校合)；法稿物理页 150,153}
\NewNote{下列是法文物理第153页的卷首异扫对照译文，属于同文校勘材料，不是新增一节正文；字迹或扫描缺损标记仍依原独立包保留。}
\subsec{超平面截面与锥形投影：拟定目录}
\begin{enumerate}
 \item 记号。
 \item 泛超平面截面：局部性质。
 \item 泛超平面截面：几何不可约性与几何连通性。
 \item 变动的超平面截面：“充分一般”的截面。
 \item 塞登伯格型定理。
 \item 任意超平面截面的连通性。
 \item 应用于构造特殊超平面截面和多截面。
 \item 例外超平面集合的维数。
 \item 改变射影浸入。
 \item 超平面截面铅笔与爆破后代数簇的纤维化。
 \item 格拉斯曼概形。
 \item 把前面的结果推广到线性截面。
 \item 初等态射与阿廷（M.~Artin）的一个定理。
 \item 锥形投影。
 \item 将前面的一些结果公理化。
 \item 改用线性系的语言表述。
\end{enumerate}

\subsec{页末说明}
这一稿将最终稿里应当收入的各项结果混在一起，作了一份详细的提要。
要写成最终稿，就得把目前这个“第零阶段”的稿子彻底重组。
第一步大概是重新拟一份计划（现第11、12、14、15条想必会大幅提前）。
我还没有写第16条。原则上它不会有任何困难，也丝毫不用借助前面各条，
因为这里要做的\gap。\cmt{L18-087-08}{本页可读到“因为这里涉及……”一类残句，后半缺损。
法第1页同样断在此处，法第2页首行接为“不过是改换表述的问题”，英第2页也有对应句。
这里只在译注中交代接文，不据此声称第153页自身完整。法第2页其余文字已见合订本编写摘要。}
\par\endgroup
